\documentclass[11pt]{article}

\usepackage[margin=1in]{geometry}
\usepackage{amsmath}
\usepackage{amssymb}
\usepackage{hyperref}
\usepackage{listings}

\lstset{
    basicstyle=\ttfamily\small,
    breaklines=true,
    columns=fullflexible,
    frame=single,
    keepspaces=true,
    showstringspaces=false
}

\title{Mathematical Engine Notes for the PRNG Casino Modules}
\author{}
\date{}

\begin{document}
\maketitle

\section{Scope}

This document explains the math behind the modules in this repository:

\begin{itemize}
    \item \texttt{prng.py}: shared deterministic HMAC-SHA256 pseudo-random generator.
    \item \texttt{game\_config.json} and \texttt{game\_config.py}: editable game math settings.
    \item \texttt{cards.py}: deck construction and unbiased Fisher-Yates shuffling.
    \item \texttt{prng\_crash\_limbo.py}: crash-game multiplier generation.
    \item \texttt{slots.py}: weighted reels, paylines, scatters, and RTP calculation.
    \item \texttt{baccarat.py}: digital baccarat dealing, totals, draw rules, and settlement.
    \item \texttt{blackjack.py}: digital blackjack dealing, hand valuation, actions, dealer rules, and settlement.
\end{itemize}

The common design is:

\[
(\text{server seed}, \text{client seed}, \text{nonce}) \longrightarrow
\text{HMAC-SHA256 output} \longrightarrow \text{game-specific outcome}.
\]

The server seed is private before the game. Its hash may be published first as a commitment:

\[
C_S = \operatorname{SHA256}(S),
\]

where \(S\) is the server seed and \(C_S\) is the public commitment. After the game, revealing \(S\) lets a player verify that the seed matches the prior commitment.

\section{Shared PRNG: \texttt{prng.py}}

\subsection{HMAC Byte Stream}

The generator creates a deterministic byte stream from:

\[
S = \text{server seed}, \qquad
C = \text{client seed}, \qquad
n = \text{nonce}, \qquad
i = \text{cursor}.
\]

For each cursor value \(i\), the message is:

\[
M_i = C \Vert \texttt{:} \Vert n \Vert \texttt{:} \Vert i.
\]

The digest block is:

\[
B_i = \operatorname{HMAC\text{-}SHA256}(S, M_i).
\]

Each digest block is 32 bytes. The byte stream is:

\[
B = B_0 \Vert B_1 \Vert B_2 \Vert \cdots.
\]

The same inputs always produce the same bytes. Without the server seed, the stream is computationally unpredictable.

\subsection{Code}

\begin{lstlisting}[language=Python]
message = f"{self.client_seed}:{self.nonce}:{self.cursor}".encode()
output += hmac.new(self.server_seed.encode(), message, hashlib.sha256).digest()
self.cursor += 1
\end{lstlisting}

The commitment helper is:

\begin{lstlisting}[language=Python]
def server_seed_hash(server_seed: str) -> str:
    return hashlib.sha256(server_seed.encode()).hexdigest()
\end{lstlisting}

\subsection{Unbiased Bounded Integers}

Many games need a random integer:

\[
X \in \{0, 1, \ldots, m - 1\}.
\]

A naive method would be:

\[
X = R \bmod m,
\]

but this creates modulo bias when the sample space size is not divisible by \(m\).

The module avoids this with rejection sampling. Let:

\[
b = \max(1, \lceil \operatorname{bitlength}(m) / 8 \rceil).
\]

The sample space is:

\[
N = 2^{8b}.
\]

The largest unbiased accepted range is:

\[
L = N - (N \bmod m).
\]

The generator draws:

\[
R \in \{0,1,\ldots,N-1\}.
\]

If \(R \geq L\), the value is rejected and another value is drawn. If \(R < L\), the result is:

\[
X = R \bmod m.
\]

Because \(L\) is divisible by \(m\), every output in \(\{0,\ldots,m-1\}\) has exactly the same number of preimages.

\subsection{Code}

\begin{lstlisting}[language=Python]
byte_length = max(1, math.ceil(upper_bound.bit_length() / 8))
sample_space = 1 << (8 * byte_length)
unbiased_limit = sample_space - (sample_space % upper_bound)

while True:
    value = int.from_bytes(self.random_bytes(byte_length), "big")
    if value < unbiased_limit:
        return value % upper_bound
\end{lstlisting}

\section{Game Configuration: \texttt{game\_config.json}}

The repository now has an editable JSON config file for target RTPs and rule settings:

\begin{lstlisting}
{
  "crash": {
    "target_rtp": 0.98,
    "instant_crash_modulus": 25
  },
  "slots": {
    "target_rtp": 0.96,
    "volatility": 1.0,
    "calibration_max_combinations": 1000000
  },
  "baccarat": {
    "target_rtp": 0.9876,
    "decks": 8,
    "banker_commission": 0.05,
    "tie_payout": 8.0
  },
  "blackjack": {
    "target_rtp": 0.995,
    "decks": 6,
    "dealer_hits_soft_17": false,
    "blackjack_payout": 1.5
  }
}
\end{lstlisting}

The values are decimals, so:

\[
0.98 = 98\% \operatorname{RTP}.
\]

For crash, \texttt{target\_rtp} is used directly as the payout scale:

\[
p = 100 \cdot \texttt{target\_rtp}.
\]

For baccarat and blackjack, RTP is primarily an output of rules and paytables. The \texttt{target\_rtp} value documents the intended target, while rule settings such as commission, tie payout, blackjack payout, dealer behavior, and deck count are what actually change the realized RTP.

For slots, \texttt{target\_rtp} can now be applied directly by calibrating a \texttt{SlotConfig}. The module first applies a volatility shape, computes the shaped exact RTP, then scales every line and scatter payout by:

\[
\alpha = \frac{\texttt{target\_rtp}}{\operatorname{RTP}_{\text{shaped}}}.
\]

After scaling, the exact RTP becomes:

\[
\operatorname{RTP}_{\text{final}}
=
\alpha \operatorname{RTP}_{\text{shaped}}
=
\texttt{target\_rtp}.
\]

This works because the slot module's payout model is linear in paytable and scatter multipliers.

Slot volatility is a numeric payout exponent:

\[
v = \texttt{volatility}.
\]

Before RTP calibration, each positive payout \(m\) becomes:

\[
m' = m^v,
\]

Values above 1 make large wins larger relative to small wins. Values below 1 compress large wins toward small wins. A value of 1 leaves payout shape unchanged. This changes payout variance while the later RTP calibration keeps the long-run average fixed.

\subsection{Code}

\begin{lstlisting}[language=Python]
def target_rtp(game: str, default: float) -> float:
    value = float(setting(game, "target_rtp", default))
    if value <= 0 or value > 1:
        raise ValueError(f"{game} target_rtp must be between 0 and 1")
    return value
\end{lstlisting}

\section{Cards and Shuffling: \texttt{cards.py}}

\subsection{Deck Construction}

For one deck:

\[
|\text{deck}| = 13 \cdot 4 = 52.
\]

For \(d\) decks:

\[
|\text{shoe}| = 52d.
\]

Cards are rank-suit pairs:

\[
\text{Card} = (\text{rank}, \text{suit}).
\]

\subsection{Fisher-Yates Shuffle}

For a deck array \(A\) of length \(N\):

\[
\text{for } i = N-1,N-2,\ldots,1:
\]

\[
j \leftarrow \operatorname{random\_int}(i+1),
\]

\[
A_i \leftrightarrow A_j.
\]

Because \(\operatorname{random\_int}\) is unbiased, every physical permutation is equally likely:

\[
\Pr(\text{specific physical order}) = \frac{1}{N!}.
\]

For \(d\) decks, duplicate rank-suit labels exist, but the underlying shoe positions are still shuffled uniformly.

\subsection{Code}

\begin{lstlisting}[language=Python]
def build_deck(decks: int = 1) -> Deck:
    return tuple(Card(rank, suit) for _ in range(decks) for suit in SUITS for rank in RANKS)

def fisher_yates_shuffle(cards: Sequence[Card], rng: FairRng) -> Deck:
    shuffled = list(cards)
    for index in range(len(shuffled) - 1, 0, -1):
        swap_index = rng.random_int(index + 1)
        shuffled[index], shuffled[swap_index] = shuffled[swap_index], shuffled[index]
    return tuple(shuffled)
\end{lstlisting}

\section{Crash Multiplier: \texttt{prng\_crash\_limbo.py}}

\subsection{Per-Game Server Seed}

The crash module derives a per-game seed:

\[
S_n = \operatorname{SHA256}(S_0 \Vert \texttt{:} \Vert n).
\]

Then it computes:

\[
D = \operatorname{HMAC\text{-}SHA256}(S_n, C \Vert \texttt{:} \Vert n \Vert \texttt{:0}).
\]

\subsection{Instant Crash}

The code checks the first 8 hex characters of the digest:

\[
Y = \operatorname{int}(D_{0:8},16).
\]

This is a 32-bit value. With modulus \(q=25\), instant crash occurs when:

\[
Y \bmod q = 0.
\]

The exact finite probability is:

\[
\Pr(\text{instant crash})
=
\frac{\left\lfloor (2^{32}-1)/25 \right\rfloor + 1}{2^{32}}.
\]

This is approximately:

\[
\Pr(\text{instant crash}) \approx \frac{1}{25} = 0.04.
\]

\subsection{Multiplier Formula}

If no instant crash occurs, the module takes 13 hex characters:

\[
h = \operatorname{int}(D_{0:13},16).
\]

Since one hex character is 4 bits:

\[
13 \cdot 4 = 52,
\]

so:

\[
0 \leq h < e, \qquad e = 2^{52}.
\]

With house edge \(r=2\%\), the payout percent is:

\[
p = 100-r = 98.
\]

The multiplier is:

\[
M = \max\left(1,\frac{pe-h}{100(e-h)}\right).
\]

As \(h\) approaches \(e\), the denominator \(e-h\) becomes small, creating rare large multipliers. Because of the \(\max(1,\cdot)\) clamp, some non-instant-crash hashes also return exactly \(1.00\). The raw expression is at most 1 when:

\[
\frac{pe-h}{100(e-h)} \leq 1.
\]

For \(p=98\), this simplifies to:

\[
h \leq \frac{2}{99}e.
\]

\subsection{Code}

\begin{lstlisting}[language=Python]
server_seed = derive_server_seed(initial_seed, game_number)
digest = hmac_sha256(server_seed, client_seed, nonce).hex()

if is_divisible(digest, INSTANT_CRASH_MODULUS):
    return 1.0

hash_int = int(digest[:13], 16)
max_hash_int = 2**52
payout_percent = 100 - HOUSE_EDGE_PERCENT

return max(1.0, (payout_percent * max_hash_int - hash_int) / (max_hash_int - hash_int) / 100)
\end{lstlisting}

\section{Slots: \texttt{slots.py}}

\subsection{Weighted Reels}

A reel is a sequence:

\[
R_k=(s_{k,0},s_{k,1},\ldots,s_{k,L_k-1}),
\]

where \(k\) is the reel index and \(L_k\) is the reel length.

The helper \texttt{expand\_weighted\_reel} converts weights into repeated symbols. If symbol \(a\) has weight \(w_a\), then:

\[
\Pr(a \text{ is selected on that reel})=
\frac{w_a}{\sum_b w_b}.
\]

\subsection{Code}

\begin{lstlisting}[language=Python]
def expand_weighted_reel(symbol_weights: Mapping[Symbol, int]) -> Reel:
    reel: list[Symbol] = []
    for symbol, weight in symbol_weights.items():
        reel.extend([symbol] * weight)
    return tuple(reel)
\end{lstlisting}

\subsection{Reel Stops}

For each reel \(k\), the PRNG chooses:

\[
x_k \in \{0,1,\ldots,L_k-1\}.
\]

The stop tuple is:

\[
\mathbf{x}=(x_0,x_1,\ldots,x_{K-1}).
\]

Because each stop is selected with unbiased \texttt{random\_int}:

\[
\Pr(\mathbf{x})=\prod_{k=0}^{K-1}\frac{1}{L_k}.
\]

\subsection{Visible Grid}

For \(r\) visible rows:

\[
G_{i,k}=R_k[(x_k+i)\bmod L_k].
\]

\subsection{Code}

\begin{lstlisting}[language=Python]
rng = FairRng(server_seed, client_seed, nonce)
stops = tuple(rng.random_int(len(reel)) for reel in config.reels)

def visible_grid(reels: Sequence[Reel], stops: Sequence[int], rows: int) -> Grid:
    columns = []
    for reel, stop in zip(reels, stops):
        columns.append(tuple(reel[(stop + row) % len(reel)] for row in range(rows)))
    return tuple(tuple(column[row] for column in columns) for row in range(rows))
\end{lstlisting}

\subsection{Payline Evaluation}

A payline is:

\[
\ell=(\ell_0,\ell_1,\ldots,\ell_{K-1}).
\]

The symbols on the line are:

\[
(G_{\ell_0,0},G_{\ell_1,1},\ldots,G_{\ell_{K-1},K-1}).
\]

For each pay symbol \(a\), the module counts consecutive symbols from the left that are either \(a\) or the wild symbol \(W\):

\[
c_a=\max c \text{ such that } G_{\ell_i,i}\in\{a,W\}
\text{ for }0\leq i<c.
\]

The selected line win is:

\[
\max_a \operatorname{paytable}(a,c_a).
\]

\subsection{Code}

\begin{lstlisting}[language=Python]
for pay_symbol, payouts in config.paytable.items():
    count = 0
    for symbol in symbols:
        if symbol == pay_symbol or symbol == config.wild_symbol:
            count += 1
        else:
            break
    multiplier = payouts.get(count)
\end{lstlisting}

\subsection{Scatter Evaluation}

If a scatter symbol \(S\) is configured:

\[
c_S=\sum_{i,k}\mathbf{1}(G_{i,k}=S).
\]

If \(\operatorname{scatter\_pays}(c_S)\) exists, it is added to the total multiplier.

\subsection{Code}

\begin{lstlisting}[language=Python]
count = sum(symbol == config.scatter_symbol for row in grid for symbol in row)
multiplier = config.scatter_pays.get(count)
\end{lstlisting}

\subsection{Total Multiplier and Payout}

\[
T=\sum_{\text{line wins } w}w+\text{scatter win}.
\]

For bet \(B\):

\[
\operatorname{payout}=BT.
\]

\subsection{Code}

\begin{lstlisting}[language=Python]
total_multiplier = sum(win.multiplier for win in line_wins)
if scatter_win is not None:
    total_multiplier += scatter_win.multiplier
payout = bet * total_multiplier
\end{lstlisting}

\subsection{Exact and Sample RTP}

Let:

\[
\Omega=\prod_{k=0}^{K-1}\{0,1,\ldots,L_k-1\}.
\]

Then:

\[
|\Omega|=\prod_{k=0}^{K-1}L_k.
\]

Exact RTP is:

\[
\operatorname{RTP}
=
\frac{1}{|\Omega|}
\sum_{\mathbf{x}\in\Omega}T(\mathbf{x}).
\]

For \(N\) sampled spins:

\[
\widehat{\operatorname{RTP}}
=
\frac{1}{N}\sum_{i=0}^{N-1}T_i.
\]

\subsection{Configurable Slot RTP}

The helper \texttt{configure\_rtp} applies the configured slot volatility and target RTP. First, it shapes all configured line and scatter payouts:

\[
\operatorname{paytable}_{\text{shaped}}(a,c) = \operatorname{paytable}(a,c)^v,
\]

\[
\operatorname{scatter\_pays}_{\text{shaped}}(c) = \operatorname{scatter\_pays}(c)^v.
\]

Then it scales all shaped payouts:

\[
\operatorname{paytable}'(a,c) = \alpha \operatorname{paytable}_{\text{shaped}}(a,c),
\]

\[
\operatorname{scatter\_pays}'(c) = \alpha \operatorname{scatter\_pays}_{\text{shaped}}(c),
\]

where:

\[
\alpha = \frac{\texttt{target\_rtp}}{\operatorname{exact\_rtp}(\texttt{shaped config})}.
\]

This should be done once when constructing a game config, not once per spin.

\subsection{Code}

\begin{lstlisting}[language=Python]
stop_ranges = [range(len(reel)) for reel in config.reels]
for stops in product(*stop_ranges):
    total_multiplier += evaluate_stops(config, stops).total_multiplier
return total_multiplier / combinations

def configure_rtp(
    config: SlotConfig,
    target: Optional[float] = None,
    volatility: Optional[float] = None,
    max_combinations: Optional[int] = None,
) -> SlotConfig:
    target = target if target is not None else target_rtp("slots", 0.96)
    shaped_config = shape_volatility(config, volatility_exponent(volatility))
    base_rtp = exact_rtp(shaped_config, max_combinations=max_combinations)
    return scale_payouts(shaped_config, target / base_rtp)
\end{lstlisting}

\section{Baccarat: \texttt{baccarat.py}}

\subsection{Card Values}

\[
v(A)=1,\qquad v(2),\ldots,v(9)=2,\ldots,9,
\]

\[
v(10)=v(J)=v(Q)=v(K)=0.
\]

The total is:

\[
H=\left(\sum_{c\in\text{hand}}v(c)\right)\bmod 10.
\]

\subsection{Code}

\begin{lstlisting}[language=Python]
def baccarat_card_value(card: Card) -> int:
    if card.rank == "A":
        return 1
    if card.rank in {"10", "J", "Q", "K"}:
        return 0
    return int(card.rank)

def hand_total(cards: tuple[Card, ...]) -> int:
    return sum(baccarat_card_value(card) for card in cards) % 10
\end{lstlisting}

\subsection{Initial Deal}

For shuffled shoe \(D\):

\[
\text{Player}=(D_0,D_2),\qquad \text{Banker}=(D_1,D_3).
\]

A natural occurs if:

\[
H_P\in\{8,9\}\quad\text{or}\quad H_B\in\{8,9\}.
\]

If there is a natural, no third cards are drawn.

\subsection{Code}

\begin{lstlisting}[language=Python]
player_hand = [deck[0], deck[2]]
banker_hand = [deck[1], deck[3]]
natural = player_total in {8, 9} or banker_total in {8, 9}
\end{lstlisting}

\subsection{Draw Rules}

If no natural occurs, the player draws if:

\[
H_P\leq 5.
\]

If the player did not draw, the banker draws if:

\[
H_B\leq 5.
\]

If the player did draw, let \(t\) be the value of the player's third card:

\[
\begin{array}{c|c}
H_B & \text{Banker draws when} \\
\hline
0,1,2 & \text{always} \\
3 & t\neq 8 \\
4 & 2\leq t\leq 7 \\
5 & 4\leq t\leq 7 \\
6 & 6\leq t\leq 7 \\
7 & \text{never}
\end{array}
\]

\subsection{Code}

\begin{lstlisting}[language=Python]
if player_third_card is None:
    return banker_total <= 5
if banker_total <= 2:
    return True
if banker_total == 3:
    return player_third_value != 8
if banker_total == 4:
    return 2 <= player_third_value <= 7
if banker_total == 5:
    return 4 <= player_third_value <= 7
if banker_total == 6:
    return 6 <= player_third_value <= 7
return False
\end{lstlisting}

\subsection{Outcome and Settlement}

\[
\text{outcome}=
\begin{cases}
\text{player},&H_P>H_B,\\
\text{banker},&H_B>H_P,\\
\text{tie},&H_P=H_B.
\end{cases}
\]

Let stake be \(B\), banker commission be \(\gamma\), and tie payout be \(\tau\). Profit is:

\[
\operatorname{profit}_{\text{player bet}}=
\begin{cases}
B,&\text{player wins},\\
0,&\text{tie},\\
-B,&\text{banker wins},
\end{cases}
\]

\[
\operatorname{profit}_{\text{banker bet}}=
\begin{cases}
B(1-\gamma),&\text{banker wins},\\
0,&\text{tie},\\
-B,&\text{player wins},
\end{cases}
\]

\[
\operatorname{profit}_{\text{tie bet}}=
\begin{cases}
B\tau,&\text{tie},\\
-B,&\text{otherwise}.
\end{cases}
\]

Returned payout is:

\[
\operatorname{payout}=B+\operatorname{profit}.
\]

\subsection{Code}

\begin{lstlisting}[language=Python]
if bet_on == "tie":
    profit = stake * tie_payout if round_result.outcome == "tie" else -stake
elif round_result.outcome == "tie":
    profit = 0.0
elif bet_on == round_result.outcome:
    multiplier = 1.0 - banker_commission if bet_on == "banker" else 1.0
    profit = stake * multiplier
else:
    profit = -stake
\end{lstlisting}

\section{Blackjack: \texttt{blackjack.py}}

\subsection{Card Values and Aces}

\[
v(2),\ldots,v(10)=2,\ldots,10,
\]

\[
v(J)=v(Q)=v(K)=10,\qquad v(A)=11\text{ initially}.
\]

If the total \(T\) is above 21 and there are aces valued as 11, each ace conversion subtracts 10:

\[
T\leftarrow T-10.
\]

The hand is soft if at least one ace remains valued as 11.

\subsection{Code}

\begin{lstlisting}[language=Python]
total = sum(blackjack_card_value(card) for card in cards)
aces = sum(card.rank == "A" for card in cards)
while total > 21 and aces:
    total -= 10
    aces -= 1
return total, aces > 0
\end{lstlisting}

\subsection{Natural Blackjack}

\[
\text{natural}\iff |\text{cards}|=2 \land T=21 \land \neg\text{from\_split}.
\]

\subsection{Code}

\begin{lstlisting}[language=Python]
def is_natural_blackjack(hand: BlackjackHand) -> bool:
    total, _ = hand_value(hand.cards)
    return len(hand.cards) == 2 and total == 21 and not hand.from_split
\end{lstlisting}

\subsection{Initial Deal}

For shuffled deck \(D\):

\[
\text{Player}=(D_0,D_2),\qquad \text{Dealer}=(D_1,D_3).
\]

\subsection{Code}

\begin{lstlisting}[language=Python]
player_hand = BlackjackHand(cards=(deck[0], deck[2]))
dealer_hand = BlackjackHand(cards=(deck[1], deck[3]))
\end{lstlisting}

\subsection{Actions}

Hit appends the next card:

\[
H'=H\Vert D_i.
\]

Stand marks the hand complete. Double down doubles the bet, draws one card, and stands:

\[
B'=2B.
\]

A split is allowed when the two cards have equal blackjack value:

\[
v(c_0)=v(c_1).
\]

The split creates:

\[
H_1=(c_0,D_i),\qquad H_2=(c_1,D_{i+1}).
\]

\subsection{Code}

\begin{lstlisting}[language=Python]
updated_hand = replace(
    hand,
    cards=hand.cards + (card,),
    bet=hand.bet * 2,
    stood=True,
    doubled=True,
)
\end{lstlisting}

\subsection{Dealer Rule}

The dealer draws while:

\[
T_D<17.
\]

If the rules say dealer hits soft 17, the dealer also draws when:

\[
T_D=17\quad\text{and}\quad\text{soft}=\text{true}.
\]

\subsection{Code}

\begin{lstlisting}[language=Python]
def should_dealer_draw(dealer: BlackjackHand, rules: BlackjackRules) -> bool:
    total, soft = hand_value(dealer.cards)
    if total < 17:
        return True
    return total == 17 and soft and rules.dealer_hits_soft_17
\end{lstlisting}

\subsection{Settlement}

Let player total be \(T_P\), dealer total be \(T_D\), hand bet be \(B\), and blackjack payout be \(\beta=1.5\) by default. Net profit is:

\[
\operatorname{profit}=
\begin{cases}
-B,&T_P>21,\\
0,&\text{player natural and dealer natural},\\
B\beta,&\text{player natural only},\\
-B,&\text{dealer natural only},\\
B,&T_D>21,\\
B,&T_P>T_D,\\
-B,&T_P<T_D,\\
0,&T_P=T_D.
\end{cases}
\]

Returned payout is:

\[
\operatorname{payout}=B+\operatorname{profit}.
\]

\subsection{Code}

\begin{lstlisting}[language=Python]
if player_total > 21:
    profit = -hand.bet
elif player_blackjack and dealer_blackjack:
    profit = 0.0
elif player_blackjack:
    profit = hand.bet * rules.blackjack_payout
elif dealer_blackjack:
    profit = -hand.bet
elif dealer_bust:
    profit = hand.bet
elif player_total > dealer_total:
    profit = hand.bet
elif player_total < dealer_total:
    profit = -hand.bet
else:
    profit = 0.0
\end{lstlisting}

\section{Expected Value and RTP}

All games can be interpreted through expected value. If outcome \(o\) has probability \(P(o)\) and payout multiplier \(m(o)\), then:

\[
\operatorname{RTP}=\mathbb{E}[m]=\sum_{o\in\Omega}P(o)m(o).
\]

House edge is:

\[
\operatorname{house\ edge}=1-\operatorname{RTP}.
\]

For slots, \(\Omega\) is the set of reel-stop combinations. For baccarat and blackjack, \(\Omega\) is the set of possible shuffled deck orders and player decisions. Baccarat has fixed decisions after the shuffle; blackjack is strategy-dependent.

\section{Math Audit Notes}

\begin{itemize}
    \item The shared PRNG uses rejection sampling, so bounded integers are unbiased as long as HMAC-SHA256 is modeled as uniformly random.
    \item Fisher-Yates is uniform because each swap index is drawn uniformly from \(\{0,\ldots,i\}\).
    \item With multiple decks, identical rank-suit labels repeat, but the underlying shoe positions are still shuffled uniformly.
    \item The crash instant-crash branch is approximately \(1/25\), with the exact finite 32-bit probability shown above.
    \item The crash multiplier can return \(1.00\) through both the explicit divisible-by-25 branch and the \(\max(1,\cdot)\) clamp.
    \item Slot RTP in this module is represented as an average payout multiplier. To express it as a percent, multiply by 100.
    \item Baccarat settlement uses net profit. A winning tie bet at \(8{:}1\) has profit \(8B\) and returned payout \(9B\).
    \item Blackjack settlement also uses net profit. A natural blackjack at \(3{:}2\) has profit \(1.5B\) and returned payout \(2.5B\).
\end{itemize}

\section{Verification Flow}

A verification-friendly round has this structure:

\begin{enumerate}
    \item Casino chooses private server seed \(S\).
    \item Casino publishes commitment \(C_S=\operatorname{SHA256}(S)\).
    \item Player chooses or accepts client seed \(C\).
    \item Game uses nonce \(n\).
    \item Outcome is computed deterministically from \((S,C,n)\).
    \item Casino reveals \(S\).
    \item Player checks \(\operatorname{SHA256}(S)=C_S\), then recomputes the game outcome.
\end{enumerate}

This proves the revealed seed matches the prior commitment. It does not prove that a paytable is favorable or that a game has a high RTP; it proves that the outcome was generated from the committed seed, client seed, and nonce.

\end{document}
